\section{Machine Learning}





\subsection {Face recognition tasks}
\subsubsection {Verification}
\subsubsection {Identification}
\subsection {Neural networks for vision}
\subsection {Embedding learning and similarity}
\section{Supervised learning}
\section{Performance metrics (verification)}
\section{Runtime metrics}
\section{Overview of models (architecture essentials)}
\section{Datasets and preprocessing}
Machine Learning (ML) has advanced rapidly in recent years, driven by increasing computational power and the availability of large-scale datasets. Despite these developments, the performance and reliability of ML models depend heavily on the quality of the data on which they are trained and evaluated. High-quality datasets are essential because they capture the relevant characteristics of real-world phenomena, provide sufficient diversity and coverage, and minimize sources of bias or noise. When datasets are incomplete, inaccurate, or unrepresentative, model performance can degrade significantly, leading to reduced accuracy, poor generalization, and unintended biases during deployment.
Understanding the dimensions of dataset quality and their influence on model behavior is therefore a crucial prerequisite for building effective and trustworthy ML systems \cite{gong2023dataquality}.